\section{ماژول bonsai\_pruner (Bonsai Pruner)}
\label{sec:module_bonsai_pruner}

\subsection{معرفی و انگیزه}
این مکانیزم بر مبنای مبانی نظری پیشرفته در ادبیات توجه و فشرده‌سازی \cite{ref_bonsaipruner_2024} توسعه یافته است.

هرس تک‌مرحله‌ای تهاجمی در شبکه‌های عصبی عمیق معمولاً به از دست رفتن پیوندهای حساس و شوک ناگهانی به بازنمایی معنایی لایه‌ها منتهی می‌شود. این آسیب در لایه‌های ابتدایی و ساختارهای چندتوجهی به ویژه بر روی توکن‌های حیاتی هدایت مانند توکن طبقه‌بندی (\lr{[CLS]}) یا سرنخ‌های توجه متمرکز در مدل‌های خودبازگشتی، بسیار برجسته‌تر است. جهت رفع این مسئله، روش‌های هرس تدریجی مکرر پیشنهاد شده‌اند، اما بدون مکانیزم‌های بازیابی بین‌مرحله‌ای، خطاهای تقریب انباشته می‌شوند.

چارچوب فشرده‌سازی \lr{Bonsai: Iterative Magnitude and Taylor Pruning with Gradient Descent Recovery} به منظور حل این نقصان ایجاد شده است. ایده بنیادین این هرس‌کننده، حذف تدریجی و چندمرحله‌ای وزن‌ها بر اساس یک زمان‌بندی مکعبی ملایم (\lr{Cubic Sparsity Schedule}) به همراه گام‌های بازیابی ریزگرادیانی (\lr{Inter-Step Gradient Descent Recovery}) میان هر دو مرحله هرس است. علاوه بر این، الگوریتم به سیستم حفاظتی مسیردهی توکن \lr{[CLS]} مجهز شده تا مهم‌ترین کانال‌های مسئول استخراج بردار جمله از اعمال هرس مصون بمانند.

\subsection{مبانی نظری و فرمول‌بندی ریاضی}
فرآیند فشرده‌سازی در \lr{Bonsai Pruner} بر پایه اصول ریاضی زیر فرمول‌بندی شده است:
\begin{enumerate}
    \item \textbf{زمان‌بندی مکعبی تنکی (\lr{Cubic Schedule}):} برای دستیابی به نسبت تنکی نهایی $S_{\text{target}}$ در طول $T$ گام متوالی، تنکی هدف در گام $t \in \{1, \dots, T\}$ طبق رابطه زیر تعیین می‌شود:
    \begin{equation}
    S_t = S_{\text{target}} \left( 1 - \left( 1 - \frac{t}{T} \right)^3 \right)
    \end{equation}
    این رابطه تضمین می‌کند که در گام‌های آغازین نرخ حذف وزن‌ها بسیار آرام و کنترل‌شده بوده و در گام‌های پایانی به آرامی به سطح هدف همگرا گردد.
    \item \textbf{معیار حساسیت تیلور مرتبه اول با حفاظت \lr{[CLS]}:} در هر گام، اهمیت هر پارامتر وزن بر اساس ضرب اندازه وزن در گرادیان تجربی سنجیده می‌شود:
    \begin{equation}
    \mathbf{S}_{ij} = |\mathbf{W}_{ij}| \cdot \left| \frac{\partial \mathcal{L}}{\partial \mathbf{W}_{ij}} \right|
    \end{equation}
    در لایه‌های توجه، سطرها و ستون‌های مرتبط با هدایت توکن \lr{[CLS]} با ضریب حفاظتی $\kappa > 1$ تقویت می‌گردند تا در اولویت حذفیات قرار نگیرند: $\mathbf{S}_{\text{CLS}} \leftarrow \kappa \cdot \mathbf{S}_{\text{CLS}}$.
    \item \textbf{گام‌های بازیابی گرادیانی بین‌مرحله‌ای (\lr{Gradient Recovery}):} پس از به‌روزرسانی ماسک تنک $\mathbf{M}_t$، وزن‌های باقیمانده به تعداد $K$ تکرار ریزگرادیانی با نرخ یادگیری کوچک $\eta_{\text{rec}}$ طبق رابطه زیر بهینه‌سازی می‌شوند تا انحراف لاجیت‌ها جبران گردد:
    \begin{equation}
    \mathbf{W} \leftarrow \mathbf{W} - \eta_{\text{rec}} \cdot \left( \nabla_{\mathbf{W}} \mathcal{L}_{\text{task}} \odot \mathbf{M}_t \right)
    \end{equation}
\end{enumerate}

\paragraph{تحلیل طیفی و نامساوی‌های پایداری خطا}
برای بررسی میزان دقت تقریب ماتریس وزن یا نگاشت توجه در ماژول \lr{bonsai\_pruner}، تفاضل ماتریس فاکتورگیری‌شده از نگاشت متراکم اصلی بر اساس نامساوی هافدینگ-فون‌نیومان ارزیابی می‌گردد. فرض کنید $\mathbf{W} \in \mathbb{R}^{d_{\text{in}} \times d_{\text{out}}}$ ماتریس وزن هدف باشد که توسط تقریب کم‌مرتبه یا پراکنده با رتبه موثر $r$ جایگزین شده است. کران بالای خطای بازسازی بر مبنای نرم فروبنیوس به صورت زیر تصریح می‌شود:
\begin{equation}
\label{eq:spectral_bound_bonsai_pruner}
\| \mathbf{W} - \mathbf{W}_{r} \|_F = \sqrt{\sum_{i=r+1}^{\min(d_{\text{in}}, d_{\text{out}})} \sigma_i^2(\mathbf{W})} \le \epsilon \| \mathbf{W} \|_F
\end{equation}
که در آن $\sigma_i(\mathbf{W})$ بیانگر $i$-امین مقدار منفرد مرتب‌شده ماتریس به صورت نزولی است. نسبت حفظ انرژی طیفی $\eta_r = \frac{\sum_{i=1}^r \sigma_i^2}{\sum_{i=1}^{d} \sigma_i^2}$ تضمین می‌کند که با انتخاب آستانه $\tau \ge 0.90$، بیش از نود درصد از پراکندگی واریانس داده‌ها حفظ شده و افت سیگنال در انتشار گرادیان رو به عقب به حداقل ممکن کاهش می‌یابد.

\paragraph{تحلیل مجانبی و پیچیدگی اشتراک حافظه سخت‌افزاری}
از منظر پیچیدگی محاسباتی و انتقال داده‌ها در سلسومراتب حافظه (حافظه سراسری \lr{HBM} به حافظه محلی \lr{SRAM})، عملیات ضرب ماتریسی در این مکانیزم نیازمند تخصیص بافر موقت با ابعاد فاکتورگیری‌شده است. برای پردازش دنباله‌ای به طول $N$ با بعد مخفی $D$، بار محاسباتی هسته از فرم درجه دو به رابطه زیر تقلیل می‌یابد:
\begin{equation}
\label{eq:complexity_bonsai_pruner}
\mathcal{C}_{\text{bonsai\_pruner}} = \mathcal{O}\left( B \cdot N \cdot D \cdot r + B \cdot H \cdot N \cdot d_k \right) \ll \mathcal{O}\left( B \cdot H \cdot N^2 \cdot d_k \right)
\end{equation}
این بهبود مرتبه پیچیدگی، با کاهش چشمگیر حجم تبادل تانسورها در گذرگاه‌های ارتباطی کارت گرافیک، مانع از اشباع پهنای باند حافظه (پدیده \lr{Memory-Bound}) در فرآیند استنتاج طولانی‌مدت می‌گردد.

\subsection{معماری و پیاده‌سازی}
کد منبع این ماژول در مسیر \lr{Codes/framework/compression\_engine/bonsai\_pruner.py} پیاده‌سازی شده و در رجیستری با شناسه‌های \lr{bonsai} و \lr{bonsai\_pruner} به ثبت رسیده است. ویژگی‌های برجسته معماری آن عبارتند از:
\begin{itemize}
    \item \textbf{کلاس \lr{BonsaiPrunedLoRALinear}:} لایه خطی اختصاصی که وزن هرس‌شده را فریز نموده و بافر ماسک را به صورت \lr{non-persistent} ذخیره می‌نماید تا حجم دیسک را کاهش دهد. این کلاس امکان تعبیه آداپتورهای سبک \lr{LoRA} را نیز فراهم می‌آورد.
    \item \textbf{کلاس \lr{BonsaiPruner}:} هدایت‌کننده حلقه اصلی هرس که با اعمال سقف مجاز بر گام‌های بازیابی (\lr{max\_recovery\_steps}) و کنترل اندازه گام‌ها (\lr{max\_step\_size})، از ناپایداری‌های عددی و واگرایی گرادیانی ممانعت به عمل می‌آورد.
\end{itemize}

\subsection{نتایج تجربی}
ارزیابی‌های تجربی ثبت‌شده برای فشرده‌ساز \lr{bonsai\_pruner} در سناریوهای مختلف در جدول \ref{tab:benchmark_bonsai_pruner} گزارش گردیده است:

\begin{table}[H]
\centering
\small
\caption{نتایج بنچ‌مارک سخت‌افزاری و زبانی فشرده‌ساز \lr{bonsai\_pruner} روی مدل‌های پایه}
\label{tab:benchmark_bonsai_pruner}
\begin{tabular}{lccccc}
\toprule
\textbf{مدل پایه \lr{(Base Model)}} & \textbf{دقت \lr{(Accuracy)}} & \textbf{شاخص \lr{F1}} & \textbf{حافظه \lr{(Peak VRAM)}} & \textbf{توان \lr{(Throughput)}} & \textbf{وضعیت} \\
\midrule
\lr{DistilBERT (SST-2)} & \lr{90.48\%} & \lr{0.9095} & \lr{357.28 MB} & \lr{139,885.85 tok/s} & \lr{PASS} \\
\lr{DistilBERT (MNLI)} & \lr{80.90\%} & \lr{0.8092} & \lr{446.34 MB} & \lr{246,370.84 tok/s} & \lr{PASS} \\
\lr{RoBERTa (SST-2)} & \lr{93.35\%} & \lr{0.9358} & \lr{503.28 MB} & \lr{14,059.48 tok/s} & \lr{PASS} \\
\lr{GPT-2 (WikiText-2)} & \lr{0.00\%} & \lr{0.0000} & \lr{816.23 MB} & \lr{1,410.97 tok/s} & \lr{PASS\_NO\_COMPRESSION} \\
\lr{Qwen} & -- & -- & -- & -- & \lr{NOT\_RUN} \\
\bottomrule
\end{tabular}
\end{table}

همان‌گونه که در جدول مشهود است، ماژول \lr{bonsai\_pruner} روی مدل‌های انکودری نتایج قابل توجه پایدار و بالایی از خود نشان داده است. در مدل \lr{RoBERTa} به دقت درخشان \lr{93.35\%} با شاخص F1 معادل \lr{0.9358} دست یافته و در مدل \lr{DistilBERT} دقت \lr{90.48\%} با توان بسیار سریع \lr{139,885} توکن بر ثانیه روی SST-2 و دقت \lr{80.90\%} با توان قابل توجه \lr{246,370} توکن بر ثانیه روی دادگان MNLI حاصل گردیده است. در مدل \lr{GPT-2} وضعیت \lr{PASS\_NO\_COMPRESSION} گزارش شده است زیرا هیچ لایه‌ای شرایط فشرده‌سازی را بدون افت خطا پاس نکرده و دست‌نخورده باقی مانده است. این ماژول بر روی مدل \lr{Qwen} اجرا نشده و در وضعیت \lr{NOT\_RUN} قرار دارد.

\subsection{محدودیت‌ها و تحلیل شکست}
بررسی‌های تحلیلی و مهندسی پیرامون عملکرد این ماژول نکات زیر را آشکار می‌سازد:
\begin{itemize}
    \item \textbf{زمان طولانی‌تر مرحله کالیبراسیون و فشرده‌سازی اولیه:} به دلیل وجود حلقه تکرارشونده $T$ گامی و اجرای چند دور ریزگرادیان بازیابی پس از هر مرحله، زمان مورد نیاز برای آماده‌سازی مدل در مقایسه با روش‌های تک‌مرحله‌ای (\lr{One-Shot}) به طور معناداری بیشتر است.
    \item \textbf{حساسیت نرخ یادگیری بازیابی:} در صورتی که نرخ یادگیری بازیابی $\eta_{\text{rec}}$ متناسب با ساختار مدل تنظیم نشود، ممکن است فرآیند بازیابی دچار نوسانات لاجیت‌ها شده و طبق آستانه \lr{max\_error} ماژول ناگزیر به بازگشت به وزن‌های چگال گردد.
\end{itemize}
